\section{Casos de Uso UML}

\subsection{Actores del Sistema}
\begin{itemize}
    \item \textbf{Atleta / Usuario Final:} Persona que ejecuta la rutina de entrenamiento, registra series y consulta las métricas de sobrecarga.
\end{itemize}

\subsection{Especificación Detallada de Caso de Uso CU-01}

\begin{table}[h!]
\centering
\caption{Caso de Uso CU-01: Registrar Serie de Entrenamiento}
\begin{tabularx}{\textwidth}{l X}
\toprule
\textbf{Caso de Uso} & \textbf{CU-01: Registrar Serie de Entrenamiento y Calcular Progresión} \\
\midrule
\textbf{Actor Principal} & Atleta / Usuario Final \\
\textbf{Precondiciones} & Existe al menos un ejercicio en el catálogo y una rutina activa. \\
\textbf{Flujo Principal} & 
1. El usuario selecciona el ejercicio actual en la sesión activa.\newline
2. El sistema muestra la sugerencia de peso y reps calculada previamente.\newline
3. El usuario ingresa la carga efectuada ($kg$), reps ejecutadas y RPE ($6.0 - 10.0$).\newline
4. El usuario presiona el botón "Completar Serie".\newline
5. El sistema persiste la serie en SQLite, inicia el cronómetro de descanso e incrementa el indicador de series.\newline
6. El motor de sobrecarga actualiza el 1RM estimado de la sesión. \\
\textbf{Flujo Alternativo} & 
\textit{3a. El usuario marca la serie como "Serie de Calentamiento":}\newline
- El sistema guarda la serie pero la excluye de los cálculos de volumen efectivo y 1RM récord.\newline
\textit{3b. El usuario indica fallo muscular (RPE 10 / RIR 0):}\newline
- El sistema ajusta el factor de fatiga acumulada para el cálculo de la siguiente sesión. \\
\textbf{Postcondiciones} & Serie guardada en SQLite; el temporizador de descanso emite notificación sonora/vibración. \\
\bottomrule
\end{tabularx}
\end{table}
